% =====================================================================
% Bab VI — Kesimpulan
% =====================================================================
\chapter{Kesimpulan}
\label{bab:kesimpulan}

Bab ini merangkum temuan utama penelitian, mengelaborasi kontribusi
keilmuan dan praktis yang dihasilkan, serta mengidentifikasi arah
pengembangan lanjutan yang bermakna. Struktur bab mengikuti tiga tujuan
penelitian yang ditetapkan pada Bab~\ref{bab:pendahuluan}: (1) merancang
dan mengevaluasi arsitektur hibrida Mamba-xLSTM-Net untuk prediksi RUL
bantalan, (2) membangun skema interpretabilitas berbasis SAE dengan pemetaan
ke frekuensi karakteristik fisik bantalan, dan (3) menganalisis konsistensi
kinerja lintas dataset dan lintas kondisi operasi.

% -----------------------------------------------------------------
\section{Kesimpulan}
\label{sec:bab6_kesimpulan}
% -----------------------------------------------------------------

Penelitian ini berhasil menjawab ketiga tujuan yang ditetapkan melalui
serangkaian eksperimen komparatif pada dua dataset \emph{run-to-failure}
publik (PHM2012 dan XJTU-SY) dengan protokol evaluasi yang ketat (tiga seed
acak, pemisahan \emph{bearing-wise}, normalisasi hanya pada data latih).

\subsection{Tujuan 1: Arsitektur Mamba-xLSTM-Net untuk Prediksi RUL}
\label{subsec:bab6_kesimpulan_t1}

Arsitektur hibrida Mamba-xLSTM-Net yang mengintegrasikan \emph{selective
state space model} (Mamba bidireksional) dan \emph{extended LSTM} (xLSTM
dengan \emph{exponential gating} dan memori matriks) melalui mekanisme
\emph{gated fusion} terbukti memberikan profil kinerja yang unggul dalam
aspek stabilitas prediksi.

Pada dataset PHM2012, Mamba-xLSTM-Net mencatat RMSE rata-rata 0,242 dengan
deviasi standar lintas \emph{seed} terendah ($\pm$0,020), menunjukkan
konsistensi prediksi tertinggi di antara ketiga arsitektur yang diuji. Pada
dataset XJTU-SY, Mamba-xLSTM-Net memperoleh PHM Score tertinggi (0,907) ---
metrik yang paling relevan dengan praktik pemeliharaan prediktif karena
memperhitungkan asimetri konsekuensi prediksi terlambat versus terlalu dini.
Hasil ini konsisten dengan hipotesis bahwa paduan kemampuan pemodelan global
jangka panjang (Mamba) dan penangkapan dinamika lokal akhir degradasi
(\emph{exponential gating} xLSTM) memberikan manfaat komplementer yang
saling memperkuat pada sinyal getaran bantalan yang bersifat non-stasioner.

Kurva pelatihan Mamba-xLSTM-Net yang terus membaik hingga \emph{epoch}
terakhir (71--75 dari 75) mengindikasikan bahwa arsitektur belum mencapai
kapasitas penuh pada anggaran pelatihan yang ditetapkan, sehingga terdapat
potensi peningkatan kinerja lebih lanjut dengan anggaran komputasi yang
lebih besar.

\subsection{Tujuan 2: Interpretabilitas SAE dan Pemetaan BPFx}
\label{subsec:bab6_kesimpulan_t2}

\emph{Sparse Autoencoder} (SAE) dengan konfigurasi \emph{Top-}$k$ ($k = 51$,
$\approx$5\% \emph{sparsity}) berhasil merekonstruksi representasi laten
Mamba-xLSTM-Net dengan \emph{loss} rekonstruksi yang sangat kecil
($5{,}12 \times 10^{-4}$ pada PHM2012; $1{,}02 \times 10^{-4}$ pada XJTU-SY),
mengkonfirmasi bahwa representasi laten model memiliki struktur yang sangat
jarang (\emph{sparse}) secara inheren.

Analisis korelasi antara aktivasi fitur SAE dan amplitudo spektrum
\emph{envelope} Hilbert pada pita frekuensi karakteristik bantalan
menghasilkan temuan yang signifikan. Pada PHM2012, \emph{hit-rate} fitur
SAE tertinggi diperoleh pada BPFI (2,3\%, atau 24 dari 1024 fitur) dengan
korelasi maksimum $r = 0{,}507$, diikuti BPFO (2,0\%). Pada XJTU-SY,
BPFO mencatat \emph{hit-rate} tertinggi (2,2\%, 23 fitur) dengan korelasi
maksimum $r = 0{,}501$. Kedua dataset menunjukkan korespondensi yang
konsisten antara fitur SAE dan frekuensi karakteristik cincin dalam
(BPFI) dan cincin luar (BPFO), yang secara fisik merupakan lokasi spalling
yang paling umum pada bantalan gelinding.

Temuan ini merupakan bukti empiris pertama bahwa model \emph{deep learning}
yang dilatih untuk prediksi RUL mempelajari representasi laten yang
berkorespondensi dengan konsep fisika klasik diagnosis getaran bantalan,
tanpa supervisi eksplisit tentang frekuensi karakteristik tersebut. Implikasi
teoritis dari temuan ini dibahas lebih lanjut pada
subbab~\ref{subsec:bab6_kontribusi_keilmuan}.

\subsection{Tujuan 3: Konsistensi Kinerja Lintas Dataset dan Kondisi}
\label{subsec:bab6_kesimpulan_t3}

Perbandingan ketiga arsitektur pada dua dataset dengan karakteristik berbeda
(PHM2012: \emph{test set} delapan bantalan, tiga kondisi operasi;
XJTU-SY: \emph{test set} dua bantalan, dua kondisi operasi) menunjukkan
beberapa pola konsistensi yang bermakna.

Pertama, peringkat kinerja antar model berfluktuasi antar dataset karena
karakteristik distribusi data yang berbeda (rekaman per bantalan, variasi
kondisi operasi, dominansi noise), namun Mamba-xLSTM-Net secara konsisten
mempertahankan \emph{trade-off} terbaik antara akurasi titik (RMSE) dan
stabilitas (deviasi standar lintas \emph{seed}) di PHM2012, serta PHM Score
terbaik di XJTU-SY.

Kedua, N-BEATS-xLSTM-RUL memberikan stabilitas \emph{seed} tertinggi
($\sigma_{\text{RMSE}} = \pm$0,003 pada XJTU-SY; $\pm$0,007 pada PHM2012),
mengindikasikan bahwa \emph{prior} struktural berbasis basis polinomial
secara efektif mengurangi sensitivitas terhadap inisialisasi acak. Sifat
ini relevan untuk \emph{deployment} produksi.

Ketiga, SparseGate-TCN-RUL menunjukkan kinerja kompetitif pada PHM2012
(RMSE 0,226, terendah) tetapi dengan variabilitas yang tinggi, mengindikasikan
bahwa pendekatan berbasis konvolusi temporal dengan \emph{sparse gating}
sensitif terhadap anggaran pelatihan yang diberikan.

% -----------------------------------------------------------------
\section{Kontribusi Penelitian}
\label{sec:bab6_kontribusi}
% -----------------------------------------------------------------

Penelitian ini menghasilkan kontribusi pada dua ranah: keilmuan teknik mesin
dan penerapan praktis dalam konteks industri Indonesia.

\subsection{Kontribusi Keilmuan}
\label{subsec:bab6_kontribusi_keilmuan}

\textbf{Kontribusi 1 (Teknologi-Metodologi)}: Arsitektur hibrida Mamba-xLSTM-Net
merupakan kombinasi pertama yang secara eksplisit memadukan \emph{selective
state space model} (Mamba bidireksional) dan \emph{extended LSTM} (xLSTM
dengan sLSTM dan mLSTM) melalui mekanisme \emph{gated fusion} yang
didesain khusus untuk karakteristik sinyal degradasi bantalan gelinding.
Penelitian yang ada (\citeauthor{Rui2026xLSTMTransformer}, \citeyear{Rui2026xLSTMTransformer})
menggabungkan xLSTM dengan Transformer (yang memiliki kompleksitas kuadratik),
sementara penelitian ini mengganti komponen Transformer dengan Mamba
(kompleksitas linear) dan memperkenalkan mekanisme fusi berbasis \emph{gate}
adaptif. Kontribusi ini relevan untuk komunitas penelitian prognostik yang
mencari arsitektur yang skalabel untuk sekuens degradasi panjang.

\textbf{Kontribusi 2 (Teknologi-Metodologi)}: Prosedur kuantitatif pemetaan
fitur SAE ke frekuensi karakteristik fisik bantalan (BPFO, BPFI, BSF, FTF)
dikembangkan dan divalidasi pada dua dataset publik. Prosedur ini terdiri
dari: (a) pelatihan SAE \emph{post-hoc} pada representasi laten model RUL,
(b) perhitungan korelasi antara aktivasi fitur SAE dan amplitudo spektrum
\emph{envelope} Hilbert pada pita BPFx, dan (c) identifikasi fitur dengan
\emph{hit-rate} dan korelasi tertinggi sebagai fitur yang berkaitan secara
fisik. Prosedur ini dapat diadopsi untuk model dan dataset lain dalam domain
prognostik mekanik.

\textbf{Kontribusi 3 (Output --- Temuan Empiris)}: Bukti empiris pertama
bahwa model \emph{deep learning} yang dilatih untuk prediksi RUL bantalan
mempelajari representasi laten yang berkorespondensi dengan frekuensi
karakteristik fisik bantalan (BPFI pada PHM2012, $r = 0{,}507$; BPFO pada
XJTU-SY, $r = 0{,}501$) tanpa supervisi eksplisit. Temuan ini mendukung
hipotesis \emph{superposition} di domain prognostik mekanik: model
mengkodekan konsep fisika domain secara implisit dalam representasi
tersembunyi yang dapat didekomposisi melalui SAE. Implikasi praktis dari
temuan ini adalah peningkatan kepercayaan rekayasawan terhadap model
\emph{deep learning} dalam sistem perawatan prediktif.

\textbf{Kontribusi 4 (Konsep-Objek + Output --- Studi Kasus Industri)}:
Studi kasus pada lini produksi PT~SKF~Indonesia (mesin \emph{grinding
outer ring} OR1 dan OR2, Pilar~0) yang mengintegrasikan sinyal getaran
spindel dengan basis data kualitas produk (\emph{vibration checking}
dan \emph{radial clearance checking}) sesuai kerangka \emph{big data}
Industri~4.0 \citetitb{Yan2017BigData}. Pada studi kasus ini dilakukan
benchmarking enam algoritma klasifikasi tiga kelas kondisi
(\emph{Normal}, \emph{Warning}, \emph{Critical}) dengan analisis
\emph{SHAP} untuk transparansi keputusan model dan validasi lintas
mesin sebagai bukti generalisasi pada konteks produksi nyata.
Temuan SHAP yang menonjolkan dominansi fitur energi pita BPFO
konsisten dengan pemetaan SAE--BPFx pada Pilar~2, sehingga ketiga
pilar disertasi (Pilar~0 industri, Pilar~1 RUL, Pilar~2 interpretabilitas)
saling memperkuat secara empiris.

\subsection{Kontribusi untuk Industri Indonesia}
\label{subsec:bab6_kontribusi_industri}

Program Making Indonesia 4.0 \citetitb{Kemenperin2018Making4} mengidentifikasi
manufaktur otomotif dan elektronik sebagai sektor prioritas transformasi
digital, dengan penerapan kecerdasan buatan dalam operasi pabrik sebagai
salah satu pilar utama. Perawatan prediktif berbasis kondisi (\emph{condition-based
maintenance}) merupakan salah satu aplikasi AI yang paling langsung dapat
diterapkan untuk meningkatkan produktivitas dan mengurangi \emph{downtime}
tidak terencana \citetitb{Mobley2002PlantMaintenance}.

Arsitektur Mamba-xLSTM-Net yang dikembangkan dalam penelitian ini dapat
diterapkan langsung pada sistem akuisisi data vibrasi yang sudah ada di
fasilitas produksi Indonesia, dengan beberapa catatan penting:

Pertama, kebutuhan komputasi pada saat inferensi relatif rendah ---
arsitektur dapat dijalankan pada CPU standar setelah pelatihan selesai,
menggunakan rekurens SSM tanpa kebutuhan GPU. Hal ini memudahkan
\emph{deployment} di lingkungan industri dengan infrastruktur komputasi
yang terbatas.

Kedua, skema interpretabilitas SAE memberikan hasil yang dapat dikomunikasikan
kepada teknisi perawatan dalam bahasa domain fisika yang akrab (frekuensi
BPFO/BPFI), bukan hanya angka prediksi abstrak. Hal ini meningkatkan
kepercayaan dan penerimaan teknologi oleh pengguna akhir.

Ketiga, arsitektur yang divalidasi pada dua dataset publik dengan kondisi
operasi yang beragam memberikan dasar yang kuat untuk adaptasi ke mesin
produksi Indonesia melalui \emph{fine-tuning} atau \emph{transfer learning}
dengan data lokal yang lebih terbatas.

% -----------------------------------------------------------------
\section{Saran Pengembangan Lanjutan}
\label{sec:bab6_saran}
% -----------------------------------------------------------------

Berdasarkan temuan dan keterbatasan yang diidentifikasi pada
Bab~\ref{bab:prognostik}, berikut ini enam arah pengembangan yang
diprioritaskan untuk penelitian lanjutan.

\textbf{Saran 1 --- Validasi dengan Anggaran Pelatihan yang Lebih Panjang}.
Kurva pelatihan Mamba-xLSTM-Net yang belum konvergen pada 75~\emph{epoch}
mengindikasikan bahwa kinerja dapat ditingkatkan lebih lanjut. Eksperimen
dengan 150--200~\emph{epoch} pada GPU dengan VRAM besar (misalnya NVIDIA
A100) diprioritaskan untuk mengkonfirmasi batas kinerja arsitektur ini.

\textbf{Saran 2 --- Uji Signifikansi Statistik Formal}.
Perbandingan kinerja antar model pada penelitian ini dilaporkan sebagai
mean~$\pm$~std atas tiga \emph{seed}. Uji signifikansi formal menggunakan
\emph{Wilcoxon signed-rank test} pada RMSE per-\emph{bearing} (yang
memerlukan minimal 5--10 \emph{run} per model) akan memperkuat klaim
keunggulan kinerja secara statistik.

\textbf{Saran 3 --- Validasi Nilai Korelasi SAE-BPFx yang Eksak}.
Nilai \emph{hit-rate} dan korelasi SAE-BPFx yang dilaporkan pada
Tabel~\ref{tab:bab5_hitrate_utama} perlu divalidasi lebih lanjut dengan
perhitungan korelasi pada set data yang lebih besar dan menggunakan
metode bootstrapping untuk mendapatkan interval kepercayaan yang andal.
Validasi kualitatif melalui visualisasi respons fitur SAE terhadap sinyal
getaran dengan fault buatan juga direkomendasikan.

\textbf{Saran 4 --- Transfer Learning ke Mesin Produksi Indonesia}.
Adaptasi model Mamba-xLSTM-Net ke data vibrasi mesin produksi nyata (misalnya
mesin \emph{grinding} di PT SKF Indonesia) melalui \emph{fine-tuning} dengan
data lokal merupakan validasi terpenting dari nilai praktis penelitian ini.
Strategi \emph{domain adaptation} yang mempertimbangkan perbedaan frekuensi
sampling, jenis sensor, dan kondisi beban antara dataset publik dan data
industri perlu dikembangkan secara sistematis.

\textbf{Saran 5 --- Deployment Edge dan Inferensi Real-Time}.
Kuantisasi model (INT8/FP16) dan kompresi melalui \emph{pruning} terstruktur
untuk \emph{deployment} pada perangkat \emph{edge} (misalnya NVIDIA Jetson
Nano) perlu dieksplorasi. Mamba, dengan kompleksitas inferensi linear dan
kebutuhan memori yang lebih kecil dari Transformer, merupakan kandidat yang
menjanjikan untuk \emph{deployment} pada perangkat komputasi terbatas di
lingkungan industri.

\textbf{Saran 6 --- Integrasi SAE End-to-End dalam Pelatihan Model}.
Pada penelitian ini, SAE dilatih secara \emph{post-hoc} setelah model RUL
selesai dilatih. Eksplorasi integrasi SAE sebagai komponen regulasi dalam
pelatihan model (penalti \emph{sparsity} pada representasi laten) dapat
menghasilkan model yang secara inheren lebih interpretatif tanpa biaya
akurasi yang signifikan. Pendekatan ini juga dapat dieksplorasi untuk
mendorong korespondensi langsung antara fitur laten dan frekuensi
karakteristik bantalan melalui supervisi fisik parsial.
